\newcommand{\C}{\mathbb{C}}
\newcommand{\Hom}{\operatorname{Hom}}
\newcommand{\GL}{\operatorname{GL}}
\newcommand{\rank}{\operatorname{rank}}

\examyear{2004}
\examera{平成16年}
\examfield{代数}
\examgroup{B}
\examnumber{1}
\examtitle{平成16年 B 第1問}
\examtag{線形代数}
\examtag{線形写像}

有限次元線形空間 \(A,B,C,D\) の次元をそれぞれ \(a,b,c,d\) とし，線形写像
\(g:B \to C\) の階数を \(r\) とする。線形写像
\[
  F:\Hom(A,B)\otimes \Hom(C,D)
    \longrightarrow \Hom(A,D)
\]
を
\[
  F(f\otimes h)=hgf
\]
で定めたとき，\(F\) の階数を求めよ。

\examyear{2004}
\examera{平成16年}
\examfield{代数}
\examgroup{B}
\examnumber{2}
\examtitle{平成16年 B 第2問}
\examtag{体論}
\examtag{群作用}
\examtag{不変体}

\(n\) を自然数とし，\(L=\C(T_1,\ldots,T_n)\) を \(n\) 変数有理関数体とする。
\(\C\) 上の \(L\) の自己同型 \(\sigma\) を
\[
  \sigma(T_i)=T_{i+1}\quad i=1,\ldots,n-1,
  \qquad
  \sigma(T_n)=T_1
\]
で定める。\(K=\{f\in L\mid \sigma(f)=f\}\) を不変部分体とする。

\begin{enumerate}
  \item 拡大次数 \([L:K]\) を求めよ。
  \item \(L\) の線形部分空間 \(\bigoplus_{i=1}^{n}\C T_i\) を \(\sigma\) の作用に関する固有空間に分解せよ。
  \item 多項式 \(f_1,\ldots,f_n \in \C[T_1,\ldots,T_n]\) で \(K=\C(f_1,\ldots,f_n)\) となるようなものを一組与えよ。
  \item \(n=6\) のとき \(L\) と \(K\) の中間体をすべて求めよ。
\end{enumerate}

\examyear{2004}
\examera{平成16年}
\examfield{代数}
\examgroup{B}
\examnumber{3}
\examtitle{平成16年 B 第3問}
\examtag{群作用}
\examtag{行列}
\examtag{軌道}

複素 \(2\times2\) 行列の全体を \(M_2(\C)\) で表す。
\[
  X=\{(A,B)\in M_2(\C)\times M_2(\C)
    \mid A^2=B^4=I,\ AB=B^3A\}
\]
とおき，\(H=\GL(2,\C)\) の \(X\) への作用を
\[
  h(A,B)=(hAh^{-1},hBh^{-1}),
  \qquad h\in H,\ (A,B)\in X
\]
で定義する。

\begin{enumerate}
  \item \(X\) の中の \(H\) 軌道の個数を求めよ。
  \item \(X\) の中の各 \(H\) 軌道は
  \(M_2(\C)\times M_2(\C)\simeq \C^8\) の閉部分集合であることを示せ。
  ただし \(\C^8\) には \(16\) 次元ユークリッド空間としての位相を入れておく。
  \item 上記の各 \(H\) 軌道に対し，その一点 \((A,B)\) の固定群
  \[
    \{h\in H\mid h(A,B)=(A,B)\}
  \]
  の複素多様体としての次元を求めよ。
\end{enumerate}

\examyear{2004}
\examera{平成16年}
\examfield{代数}
\examgroup{B}
\examnumber{4}
\examtitle{平成16年 B 第4問}
\examtag{可換環}
\examtag{単元群}
\examtag{有限群}

\(p\) を \(3\) 以上の素数とする。単位元をもつ可換環 \(A\) の可逆元のなす群 \(A^\times\) は，
位数 \(p^2\) の巡回群にはならないことを証明せよ。
